% 07-storage.tex
\chapter{Storage}
\label{ch:storage}

\section{Storage Drivers}
\label{sec:storage-drivers}

Doki supports 5 storage drivers, auto-detected by
\texttt{DetectBestDriver()}. The driver selection depends on the
kernel version, filesystem type, and whether root access is available.

\begin{table}[htbp]
  \centering
  \caption{Storage driver comparison.}
  \label{tab:storage-drivers}
  \begin{tabularx}{\textwidth}{l l l l X}
    \toprule
    \textbf{Driver} & \textbf{Root?} & \textbf{Performance} & \textbf{Status} & \textbf{Use case} \\
    \midrule
    overlay2       & Yes & Best   & Tested & Linux servers with kernel overlay \\
    fuse-overlayfs & No  & 90\%   & Tested & Rootless, Termux, Android \\
    aufs           & Yes & Good   & Tested & Legacy Linux systems \\
    vfs            & No  & Slow   & Tested & Fallback, debugging \\
    devmapper      & Yes & Good   & Tested & Enterprise, thin provisioning \\
    \bottomrule
  \end{tabularx}
\end{table}

\section{Overlay Filesystem}
\label{sec:overlay}

The overlay filesystem is the primary storage mechanism for container
images. It combines multiple read-only layers into a single unified
view using copy-on-write semantics.

\subsection{Layer Structure}
\label{sec:layer-structure}

An OCI image consists of multiple layers, each representing a set of
filesystem changes. When a container is created, a thin writable layer
(\texttt{upperdir}) is placed on top of the read-only image layers.

\begin{figure}[htbp]
  \centering
  \begin{tikzpicture}[
    node distance=0.3cm,
    every node/.style={font=\sffamily\footnotesize},
  ]
    % Layers (bottom to top)
    \node[rectangle, draw=nodeblue!50, fill=nodeblue!8,
      minimum width=6cm, minimum height=0.8cm,
      text=textdark, align=center, inner sep=3pt]
      (base) at (0,0) {Base Layer (alpine:3.18)};

    \node[rectangle, draw=nodelteal!50, fill=nodelteal!8,
      minimum width=6cm, minimum height=0.8cm,
      text=textdark, align=center, inner sep=3pt]
      (layer1) at (0,1.1) {Layer 1: apt-get install dependencies};

    \node[rectangle, draw=nodeorange!50, fill=nodeorange!8,
      minimum width=6cm, minimum height=0.8cm,
      text=textdark, align=center, inner sep=3pt]
      (layer2) at (0,2.2) {Layer 2: COPY app /opt/app};

    \node[rectangle, draw=nodepurple!50, fill=nodepurple!8,
      minimum width=6cm, minimum height=0.8cm,
      text=textdark, align=center, inner sep=3pt]
      (layer3) at (0,3.3) {Layer 3: RUN go build -o /app};

    \node[rectangle, draw=nodecoral!50, fill=nodecoral!12,
      minimum width=6cm, minimum height=0.8cm,
      text=textdark, align=center, inner sep=3pt]
      (writable) at (0,4.4) {Writable Layer (container)};

    % Annotations
    \node[annot, right=0.3cm of base.east] {read-only};
    \node[annot, right=0.3cm of layer1.east] {read-only};
    \node[annot, right=0.3cm of layer2.east] {read-only};
    \node[annot, right=0.3cm of layer3.east] {read-only};
    \node[annotBold, right=0.3cm of writable.east] {read-write};

    % Bracket
    \draw[decorate, decoration={brace, amplitude=5pt, mirror},
      linegray, thick]
      ([xshift=-0.2cm]base.north west) -- ([xshift=-0.2cm]writable.south west)
      node[midway, left=8pt, font=\sffamily\tiny, text=textgray, rotate=90]
        {Image Layers};

    % Merge arrow
    \draw[arrThick, nodeblue] ([xshift=3.5cm]writable.north east)
      -- ++(0.8,0) node[right, font=\sffamily\footnotesize, text=nodeblue]
        {merged view};
  \end{tikzpicture}
  \caption{Overlay filesystem layer structure. Multiple read-only
    image layers are merged with a writable container layer.}
  \label{fig:overlay-layers}
\end{figure}

\section{Image Layer Management}
\label{sec:layer-management}

Layers are stored as content-addressable blobs identified by their
SHA-256 digest. The \texttt{ImageStore} preserves full layer IDs
without truncation, while the CLI truncates IDs for display purposes.

\subsection{Layer Deduplication}
\label{sec:layer-dedup}

When pulling images, Doki tracks already-pulled layers in a map to
avoid duplicate downloads. If a layer already exists in the local
store, it is skipped during the pull process.

\section{Volume Management}
\label{sec:volume-management}

Doki supports named volumes that persist beyond container lifecycle.
Volumes are stored under \texttt{\$DOKI\_DATA\_DIR/volumes/} and can
be mounted into containers at any path.

\subsection{Volume Types}
\label{sec:volume-types}

Doki supports three volume types:

\begin{description}[leftmargin=*, style=nextline, font=\sffamily\bfseries]
  \item[Named volumes.]
    Managed by Doki with automatic lifecycle management. Data persists
    across container restarts and removals.

  \item[Bind mounts.]
    Direct mapping of host directories into the container. Useful for
    development where files need to be shared between host and
    container.

  \item[tmpfs mounts.]
    In-memory filesystems for sensitive data that should not persist
    to disk. Limited by available RAM.
\end{description}

\begin{codebox}[title={Volume mount syntax}]{bash}
doki run -v mydata:/data alpine:3.18 sh
doki run -v /host/path:/container/path alpine:3.18 sh
doki run --tmpfs /tmp:rw,size=100m alpine:3.18 sh
\end{codebox}

\section{Filesystem Probes}
\label{sec:filesystem-probes}

Doki detects its execution environment by probing the filesystem
rather than relying on \texttt{runtime.GOOS}, which returns
\texttt{linux} on both Android and standard Linux:

\begin{codebox}[title={Environment detection logic}]{go}
func IsTermux() bool \{
    \_, err := os.Stat("/data/data/com.termux")
    return err == nil
\}
\end{codebox}

\subsection{Detection Probes}
\label{sec:detection-probes}

The following filesystem paths are probed to determine the execution
environment:

\begin{table}[htbp]
  \centering
  \caption{Filesystem probes and their detected environments.}
  \label{tab:fs-probes}
  \begin{tabularx}{\textwidth}{l X}
    \toprule
    \textbf{Path} & \textbf{Environment} \\
    \midrule
    \texttt{/data/data/com.termux} & Android/Termux \\
    \texttt{/proc/self/status}     & Standard Linux (always present) \\
    \texttt{/proc/1/cgroup}        & Container detection \\
    \texttt{/.dockerenv}           & Docker container \\
    \texttt{/run/.containerenv}    & Podman container \\
    \bottomrule
  \end{tabularx}
\end{table}

This approach enables environment-specific behavior (such as DNS port
selection and proot optimization) without compile-time restrictions.
